\begin{frame}{자주 하는 실수: 학습 곡선 한 번만 보고 결론 내리기}
  \small
  \begin{itemize}
    \item RL 프로젝트에서 특히 빈번한 부실한 리뷰: ``학습 곡선이 잘
      올라갔으니 성공''로만 판단 — 그 곡선이 \textbf{무작위 시드 하나}인지,
      \textbf{반복해서 일관된 경향}인지, \textbf{완전 수렴}인지
      ``개선됐지만 갈 길이 남은'' 상태인지를 \textbf{먼저} 물어야 한다.
    \item 신경망 RL에서 \emph{구조적으로} 더 위험한 이유: 곡선의
      ``오르다/안 오르다''가 \textbf{학습의 진전과 분리될 수 있기}
      때문 — 11.2에서 PPO 곡선은 ``완만한 상승''이었으나 개별
      에피소드 리턴은 $-1800$(쓰러짐) $\sim 0$(수직)까지 넓게
      퍼져 있음. 이동평균이 올라간다고 \textbf{에피소드 리턴의 분산}
      이 줄어든 것은 아니다.
  \end{itemize}
  \vskip0.4em
  \begin{block}{학습 곡선 3문 — 리뷰어가 던지는 질문}
    \begin{enumerate}
      \item \textbf{시드(초기화)}: 시드 \emph{몇 개}의 평균(또는 범위)
        인가? 1개라면 ``일관된 경향''을 말할 수 없다.
      \item \textbf{평가 프로토콜}: ``학습 도중 리턴''(탐험 포함)인지
        ``학습 후 탐험 0 재평가 리턴''인지? \textbf{다른 대상}의
        측정값이다.
      \item \textbf{수렴 여부}: 이 스텝 수(예: 12만)가 완전
        수렴인가? 더 오래 돌리면 포화인가요, 계속 오르나요(11.2)?
    \end{enumerate}
  \end{block}
\end{frame}
